\renewcommand{\stagefigureprefix}{2}\setcounter{figure}{0}
\reporttitle{第二阶段：指令设计与基础测试}

\section{设计目标与实施范围}
点积把两组数据逐项相乘，再将乘积求和，是量化矩阵运算的重要计算步骤。本设计在昆明湖V2已有向量路径中增加八项INT8点积指令，将乘法与归约组合为专用数据通路，同时复用处理器的译码、寄存器重命名、调度、写回和精确提交机制。

本次提交和评估的是弹性两槽实现：乘积槽P与结果槽R组成两级流水，接收请求后最早2拍给出结果，每拍可以接收一笔新请求。工程同时保留更早的单请求实现，可在构建时切换，用于对照和回退。本章依次介绍指令语义、两槽数据通路与控制、处理器接入和三类基础测试，分别对应赛题2.1的设计说明、2.2的RTL实现和2.3的三类基础测试；十六项点积和INT16参数化两项扩展列在本章末尾。主要结果是：指令按赛题语义实现，正数、负数和边界值三类测试的完整128位结果均与期望一致；两槽实现达到每拍发起一笔、最早2拍响应。

图\ref{fig:S2-11}给出vdot在昆明湖V2流水线中的位置。vdot与普通向量指令一样经过取指、译码、重命名和派发；译码器按custom-0编码识别它，派发到向量执行单元VFEX4中新增的VDot功能单元，结果经共享写回送入向量寄存器堆，最后由ROB按程序顺序提交。新增和修改的部分只有译码识别、VDot功能单元和VFEX4的派发与写回选择，其余流水线机制全部复用。

\reportfigure{S2-11}{vdot在昆明湖V2流水线中的位置}{7cm}{取指、译码、重命名、派发、执行、写回和提交七段；新增与修改部分为custom-0译码、VFEX4中的VDot功能单元及共享派发与写回选择。}{依据本章译码、派发与VFEX4接入说明绘制的简化流水线示意。}

表\ref{tab:s2-decisions}汇总了本设计的几项主要决策。

\begin{table}[!htbp]\centering
\small
\caption{主要设计决策}\label{tab:s2-decisions}
\begin{tabular}{>{\raggedright\arraybackslash}p{0.13\linewidth}>{\raggedright\arraybackslash}p{0.30\linewidth}>{\raggedright\arraybackslash}p{0.48\linewidth}}
\toprule
决策 & 做法 & 理由与效果\\
\midrule
指令编码 & custom-0编码，译码为单个微操作 & 使用RISC-V为自定义扩展保留的编码空间，不占用标准RVV编码；复用现有译码、重命名和精确异常机制\\
接入位置 & 放入已有的向量执行单元VFEX4，与整数DIV、浮点DIV/SQRT共享派发和写回 & 不新增发射端口和写回端口；需要处理与DIV的互斥、取消和时钟门控（见第\ref{subsec:2-vfex4}节）\\
流水结构 & 弹性两槽：乘积槽P＋结果槽R & 相邻请求的乘法和归约重叠执行，发起间隔从上一版单笔在途设计的2拍缩短到1拍；独立批量周期减少39\%（见第三阶段）\\
结果保存 & R槽只保存32位，输出时补96位零 & 利用结果高96位恒为零的性质减少寄存器，写回接口仍是完整的128位\\
运算拍数 & 乘法和归约分两拍完成 & 另实现并验收了一拍方案SP1；两拍方案面积更小、VDot路径时序余量更大，综合权衡后采用（见第\ref{subsec:2-tradeoff}节）\\
DIV取消修复 & DIV请求被取消后不再保持时钟 & 取消后VFEX4可以完全关钟，测量周期不变\\
\bottomrule
\end{tabular}
\end{table}

\section{指令语义与编码}
\subsection{八项有符号点积与结果布局}
本指令采用128位向量寄存器（VLEN=128）、8位元素（SEW=8）和单寄存器分组（LMUL=1）；有效元素数为8（VL=8），起始执行索引为0（vstart=0）。指令读取两个源寄存器的低8个字节，将每个字节解释为有符号整数，计算
\[
S=\sum_{i=0}^{7}\mathrm{int8}(vs1[8i+7:8i])\,
                    \mathrm{int8}(vs2[8i+7:8i]).
\]
目的低32位保存$S$的二补码，高96位写零。运算从零开始累加，读取指令执行前的源值，因此支持目的与源寄存器重合，也支持使用v0作为普通源或目的。每个源以128位向量传入数据通路，八项算术使用其中低64位的8个元素。

图\ref{fig:S2-01}汇总了指令字段、源操作数取位和结果布局。custom-0编码给出两个源寄存器和一个目的寄存器；两源各取低64位的8个有符号字节，同编号字节两两相乘后求和；目的低32位保存和的补码，高96位写零。以$S=-8$为例，完整结果为\seqsplit{0x00000000\_00000000\_00000000\_FFFFFFF8}。

\reportfigure{S2-01}{八项点积指令字段与完整结果布局}{9cm}{custom-0编码选择源、目的寄存器；两源各取低64位的8个有符号字节求点积。目的低32位保存结果补码，高96位写零；负数示例同样遵循这一完整128位布局。}{八项点积指令规格与实际译码字段的设计示意。}

\subsection{数值范围与中间位宽}
有符号8位数的范围为$-128$至127，单个乘积范围为$[-16256,16384]$，八项和范围为$[-130048,131072]$。乘积使用16位；按二项、四项、八项归并时，精确位宽依次增加为17、18、19位。正端点131072比18位有符号数的最大值大1，因此极值用例能直接检验归约过程中是否保留了完整进位。

归约结果先符号扩展为32位，再与96位零拼接。负数的符号由低32位字段表达，例如$-8$的低32位为\texttt{0xfffffff8}，完整128位结果仍以96个零作为高位。也就是说，19位用于保证求和精确，32位是结果槽实际保存的宽度，128位是写回向量寄存器的架构宽度。

\subsection{编码与合法执行状态}
指令采用custom-0编码，\texttt{MATCH=0x0200000b}、\texttt{MASK=0xfe00707f}，各字段取值见表\ref{tab:s2-encoding}。
\begin{table}[!htbp]\centering
\small
\caption{vdot指令的编码字段}\label{tab:s2-encoding}
\begin{tabular}{lll}
\toprule
字段 & 位段 & 含义或取值\\
\midrule
funct6 & 31:26 & \texttt{000000}\\
vm & 25 & \texttt{1}\\
vs2、vs1 & 24:20、19:15 & 两个向量源寄存器编号\\
funct3 & 14:12 & \texttt{000}\\
vd & 11:7 & 向量目的寄存器编号\\
opcode & 6:0 & custom-0，\texttt{0x0b}\\
\bottomrule
\end{tabular}
\end{table}
汇编记法为\texttt{vdot.vv vd, vs2, vs1}。

例如\texttt{vdot.vv v10,v8,v9}对应\texttt{0x0284850b}。标准汇编器不认识这条自定义指令，程序通过经过校验的编码宏或直接写入机器字来生成它。

向量功能启用、vill=0、vm=1以及本节的SEW、LMUL、VL和vstart共同确定合法状态，四种vta/vma策略组合均适用。合法执行完成一次完整目的写回，并保持VL、vtype、vlenb、vxrm、vxsat和vcsr。非法请求携带illegal-instruction异常完成，同时关闭向量与v0目的写使能，由处理器的精确异常机制处理。

\section{数据通路与弹性两槽控制}\label{sec:2-datapath}
\subsection{乘积槽与结果槽}
点积数据通路包含8个有符号$8\times8$乘法器、乘积槽P、归约网络和结果槽R。P保存8个16位乘积及这笔请求的控制、性能信息；R保存32位结果及同一请求的身份。结果从R输出时再补入96位零，这样R只需保存32位，写回接口仍是完整的128位。

图\ref{fig:S2-02}给出这条数据通路的结构。8个乘法器的16位乘积存入P，组合归约网络按4个17位、2个18位、1个19位逐级相加得到精确和，符号扩展为32位后存入R，输出时再与96位零拼接成128位结果。图下部为事务身份通路：ROB、目的寄存器、写使能、异常信息和性能字段存入pCtrl与pPerf，P转入R时整包移入rPacket；若P中已锁存非法指令标记，选择器以零代替结果，并关闭目的写使能。

\reportfigure{S2-02}{八项点积的乘积槽、归约网络和结果槽}{8cm}{P保存8路16位乘积，组合归约依次保留17、18、19位，并符号扩展为32位后进入R。控制与目的身份随数据移动；R中的32位结果在输出时补入高96位零，异常请求按图示路径处理。}{固定八项两槽FU及32位结果保存结构的设计示意。}

两槽的有效位表示各自是否保存一笔请求。重排序缓冲区（ROB）编号、目的物理寄存器、向量或v0写使能、异常信息以及性能字段与数据同时移动，使写回始终能够找到正确的目的。输出停顿时，R中已有的请求保持不变，不受新输入影响。

\subsection{时钟、复位与连续流动}
VDot功能单元（FU）接入共享执行单元VFEX4，后文将VFEX4称为父级。实现沿用功能单元所在的单时钟域和父级复位连接。P、R在时钟上升沿更新；复位清除两槽有效位，并抑制复位期间的输入和输出握手。在父级选择并允许接收VDot、独立请求连续供给、下游持续就绪且没有复位或重定向时，请求于E0进入P，E1完成归约并进入R，E2是最早的输出握手沿。连续请求的发起间隔为1拍。表\ref{tab:s2-slot-actions}列出P、R与下游状态的典型组合及对应的本拍动作。

\begin{table}[!htbp]\centering
\small
\caption{两槽状态与本拍动作}\label{tab:s2-slot-actions}
\begin{tabular}{p{0.16\linewidth}p{0.16\linewidth}p{0.16\linewidth}p{0.40\linewidth}}
\toprule
P状态 & R状态 & 下游状态 & 本拍动作\\
\midrule
空 & 空 & 任意 & 接受新请求到P\\
有效 & 空 & 任意 & 原P进入R，同时可接收下一请求\\
空 & 有效 & 反压 & 保持R，P可接收一笔请求\\
有效 & 有效 & 反压 & 保持两槽，暂停接收\\
有效 & 有效 & 就绪 & 原R输出、原P进入R、新请求进入P\\
\bottomrule
\end{tabular}
\end{table}

图\ref{fig:S2-05}以四笔连续请求为例展示逐拍流动，图中P、R两行为时钟沿前的状态。下游持续就绪时，\#1在cycle 10被接收，cycle 11位于P，cycle 12位于R并完成响应；在cycle 12这一拍，接收\#3、\#2由P迁入R、\#1从R响应三件事同时完成。由此得到1拍的发起间隔和接收后2拍的最早响应。

\reportfigure{S2-05}{连续请求在两槽中的逐拍流动}{8.0cm}{持续供给且下游就绪时，旧R输出、旧P推进和新请求接收可以在同一拍重叠；最早响应为接受后2拍，发起间隔为1拍。}{逐拍机制示意，依据弹性槽状态方程和连续请求验证。}

全核运行的观察器日志也记录到了这种同拍重叠，如图\ref{fig:SUP-S2-OVERLAP}所示：10986周期同时记录了ACCEPT seq2、MOVE seq1和RESPONSE seq0，分别对应新请求进入P、原P推进至R和原R响应。

\evidencefigure{SUP-S2-OVERLAP}{连续事务的同拍重叠记录}{11cm}{10986周期同时记录ACCEPT seq2、MOVE seq1与RESPONSE seq0，分别对应新输入、原P推进和原R响应。}

FU能否接收取决于自身是否有空槽，但新请求何时派发仍由父级根据当前选择和共享资源占用决定。

时延指一笔请求从被接受到产生响应所用的拍数，发起间隔指相邻两笔请求被接受的间隔。两槽结构支持连续独立点积；程序的最终周期还由数据依赖、调度、访存和写回共同决定。FU通过valid/ready报告完成，下游反压时响应相应顺延。

\subsection{背压、重定向与事务守恒}
R有效而下游未就绪时，数据、目的和控制包一起保持。P可利用空槽再接收一笔请求；若两槽均占用且R仍受阻，输入ready撤销，直到R输出或取消腾出空间。

图\ref{fig:S2-06}给出一次实际背压过程，A、B分别为序号1321和1322的请求。A在+2拍进入R，下游从+2到+8连续7拍未就绪，期间R保持A的结果与目的身份，P中的B也保持不变；+9拍下游恢复就绪，A完成一次响应，B同拍由P迁入R。此时FU已经腾出空槽（FU的in.ready为1），但父级的in.ready仍为0：FU有空槽并不等于立即接收新请求，何时派发仍由父级决定。

\reportfigure{S2-06}{背压期间的数据与目的身份保持}{8.0cm}{R在输出阻塞期间保存完整响应；空闲P可继续接收一笔请求，满槽后向上游施加反压。}{控制机制示意，依据FU和真实父级反压验证。}

重定向按两个槽各自的ROB年龄独立判断取消，保留较老的幸存请求，并在该拍暂停新请求接收。被取消的R停止输出，被取消的P释放待归约状态；复位的优先级最高。ROB环回场景使用完整指针年龄关系确定先后。

图\ref{fig:S2-07}给出四种取消与复位场景，均假设下游未就绪（out.ready=0），且重定向当拍不接收新请求。只取消P时，R中的B保留，继续等待下游就绪；只取消R时，P中的A推进到R；两槽同时被取消时都变为无效；reset与redirect同拍时以reset为准，两槽清空。

\reportfigure{S2-07}{部分取消、全部取消与复位恢复}{8.0cm}{分别判断P、R中的请求是否被重定向命中，展示幸存请求、取消清空和复位后重新接收。}{控制机制示意，依据槽内ROB判断和对应定向用例。}

每次验证在统一时钟坐标下检查
\[
N_{\mathrm{accept}}=N_{\mathrm{response}}+N_{\mathrm{cancel}}
                  +N_{\mathrm{reset}}+N_{\mathrm{inflight}},
\qquad 0\leq N_{\mathrm{inflight}}\leq2.
\]
该关系将每笔已接受请求归入响应、取消、复位丢弃或当前占用之一，并在运行结束时检查两槽排空。每笔响应还核对完整128位结果、ROB和目的身份。

图\ref{fig:S2-10}综合了正常推进与取消、复位处理。A部分以(P.valid, R.valid)表示状态，标出没有flush和reset时各状态之间的转移及触发事件；B部分给出取消与复位的判定顺序：reset优先清零两槽，未复位时按各自ROB匹配取消，幸存请求仍可移动或输出，flush有效期间禁止新请求进入；C部分说明接收、移动和输出三件事可以在同一边沿重叠。

\reportfigure{S2-10}{两槽同拍推进与取消、复位处理}{10.5cm}{状态按(P.valid,R.valid)排列。满槽时原R输出、原P推进与新请求接受可在同一边沿重叠；复位优先清空，取消按各自ROB匹配，幸存请求仍可推进或输出，flush有效时禁止新请求进入。}{实际有效位更新、取消优先级及事务守恒关系的状态示意。}

\section{RTL实现与处理器接入}
\subsection{关键实现语句}
算术部分位于\seqsplit{src/main/scala/xiangshan/backend/fu/vector/VDotDatapath.scala}，乘积和每级归约按有符号类型计算，并逐级保留进位，对应的源码摘录见图\ref{fig:SUP-S2-ARITH}。\seqsplit{wrapper/VDotElastic.scala}保存P、R的控制和数据，图\ref{fig:SUP-S2-STATE}依次摘录了其中的32位结果槽寄存器、握手与取消条件，以及状态更新和原P内容转入R的逻辑。

\evidencefigure{SUP-S2-ARITH}{八项点积算术实现摘录}{12cm}{八路有符号乘法和逐级归约，对应本章的中间位宽分析。}
\evidencefigure{SUP-S2-STATE}{两槽状态与控制实现摘录}{13cm}{结果槽保存32位；有效位、握手、取消条件和状态更新共同管理请求。}

下面摘示从P转入R的关键语句：
\begin{Verbatim}[fontsize=\small,breaklines=true]
when(move) {
  rPacket.ctrl := pCtrl
  rPacket.perfDebugInfo := pPerf
  rPacket.res.data := Mux(pCtrl.exceptionVec.get(ExceptionNO.illegalInstr),
    0.U(32.W), reduce.io.result(31, 0))
}
\end{Verbatim}
这些赋值读取本拍边沿前的P内容，因此同拍接收的新请求会进入P，原请求的结果和身份则进入R。公开输出使用：
\begin{Verbatim}[fontsize=\small,breaklines=true]
io.out.bits.res.data := Cat(0.U(96.W), rPacket.res.data)
\end{Verbatim}
结果寄存器的保存位宽为32位，高96位由输出组合逻辑补零。

\subsection{译码、配置依赖与写回}
精确编码匹配后，译码器将vdot归入向量运算，生成单个微操作并建立两个向量源依赖。实际VL通过重命名后的依赖传递，SEW、LMUL和其他向量控制随指令保存。这样，多条向量配置与运算交错时，每条点积仍使用与自己对应的配置。

旧目的值不参与累加，译码去除相应的旧目的依赖；向量和v0的目的分配、异常信息及最终写回沿用现有处理器机制。两槽实现通过构建开关\texttt{VDOT\_O2\_ELASTIC=1}和\texttt{VDOT\_T14\_FUSED\_INPUT=1}启用；不加开关时构建单请求实现，便于对照和回退。

\subsection{共享VFEX4的派发与门控}\label{subsec:2-vfex4}
VDot与整数DIV、浮点DIV/SQRT共用VFEX4父级。父级根据选择、实际占用和被选功能单元的ready进行原子派发，使父级和子单元接受的是同一笔请求。VDot两槽排空后可以切换到DIV；DIV在途期间保存自己的ROB身份和占用状态。输出按唯一有效完成来源选择共享写回。

图\ref{fig:S2-03}是这一共享结构的简化示意。派发与互斥逻辑按选中的FU接收请求，三个功能单元各有独立的时钟门控，完成结果经共享写回从唯一有效来源输出；占用与控制逻辑汇总各单元的状态，约束派发，并产生取消和时钟使能。图中省略了部分逐信号连线。

\reportfigure{S2-03}{VFEX4中的派发、资源互斥与共享写回}{8cm}{固定八项两槽组合的简化结构示意。VDot、整数DIV和浮点DIV/SQRT共享派发与写回；图中省略部分逐信号连线，接收许可、ready选择、占用与门控关系按下文说明。}{依据固定受测父级派发与写回实现，简化图不作为完整接口接线图。}

派发许可同时检查资源占用、请求选择以及flush和reset。上游ready由许可与被选FU的ready共同决定，输入valid也仅送往被选单元，因此父级接受与该FU接受必须发生在同一边沿。ready的返回路径按选择信号取值，不能将各FU的ready直接作或运算。

VDot占用由\texttt{pValid}或\texttt{rValid}表示；DIV侧还保留在途ROB和占用状态，接受后持续保持，直到输出握手或匹配取消。图\ref{fig:S2-03}中的输出ready向各FU广播，但输出valid最多一路有效，共享写回按该唯一有效来源同时选择结果与身份。

时钟使能覆盖存活请求、当前输入接受和必要的取消、复位清理。对应验证包含满槽背压、部分取消、ROB环回、关钟后的复位恢复，以及VDot和两类DIV之间的切换。DIV取消修复让被取消的DIV请求不再占住时钟，父级因此能进入空闲并接受后续工作。

图\ref{fig:S2-08}给出两段实测过程。上半部分是VDot排空后切换到DIV：旧R输出、VDot排空和DIV接收依次相隔1拍，整数DIV与浮点DIV从接收到响应分别用10拍和8拍。下半部分是取消清理后的时钟变化：DIV在+0拍接收、+1拍完成取消清理，+2至+4拍父级全部关钟，+5拍重新接收VDot，+7拍给出响应。也就是说，取消清理完成后DIV不再保持时钟，父级可以完全关钟，之后再正常切换回VDot。

\reportfigure{S2-08}{共享资源切换、取消与时钟保持}{8.0cm}{按存活请求保持时钟，在取消或完成后释放相应占用，随后恢复派发。}{门控机制示意，依据真实父级的取消、空闲和恢复检查。}

\section{基础功能与分层验证}\label{sec:2-verify}
验证首先用独立数学计算确定完整128位结果，再依次检查算术RTL、带握手的功能单元、真实父级和全核执行。算术测试关注有符号乘法、极值和逐级位宽；FU与父级测试加入背压、取消、复位、目的身份和DIV交互；整核协同验证进一步检查译码、架构状态及提交结果。

图\ref{fig:S2-04}按测试输入、五层验证和期望与检查三列组织这些测试。算术测试输入同时驱动独立数学模型和算术RTL，前者给出完整128位期望，与后者的实际值比较；分层事务用例分别驱动VDot FU和真实VFEX4父级，检查反压、取消、复位、身份，以及父级的互斥、切换和空闲关钟；AM功能矩阵在全核EMU与NEMU上运行，通过DiffTest和程序自检核对架构状态。算术RTL的10022项输入与全核117个独立用例、124次执行分别统计，各层证据独立归档。

\reportfigure{S2-04}{从数学结果到整核执行的验证层次}{9.5cm}{数学模型提供完整128位期望值，算术RTL、FU、真实VFEX4父级和全核分别核对数值、握手与架构状态。算术10022项输入与全核117个独立用例、124次执行分别统计，各层证据独立归档。}{已有测试入口与结果记录的分层验证示意；测试统计不跨层合并。}

以下三类测试对应赛题2.3.1至2.3.3。图\ref{fig:S2-TEST-01}至图\ref{fig:S2-TEST-03}的输入与响应取自固定八项两槽版本的全核观察记录，按seq、ROB和目的物理寄存器配对；第三阶段说明该版本的架构差分验证。
\clearpage
\subsection{题目2.3.1：正数点积测试}
两组有效输入均为[1,2,3,4,5,6,7,8]，八项平方和为204。历史全核functional-2记录中，seq=1、ROB=0:39、pdest=103在339296周期接受、339298周期响应；实际完整128位结果为\seqsplit{0x000000000000000000000000000000cc}，与预期相同。源高64位为A5哨兵，不参与点积。下图给出对应输入与实际响应摘录。\par
\reportfigure{S2-TEST-01}{正数点积测试：预期与实际结果均为204}{5.0cm}{两个源的有效元素均为1至8；expected与actual完整128位一致，末尾为CC。}{用户保存的原始PNG；functional-2观察器日志字段摘录，保留事务身份、原始行号与SHA256；不是波形截图。}
\subsection{题目2.3.2：负数点积测试}
有效输入为八个−1与八个1，预期总和−8。functional-0的seq=4、ROB=0:12、pdest=143在420291周期接受、420293周期响应；实际完整结果为\seqsplit{0x000000000000000000000000fffffff8}。低32位保存负数补码，高96位为零；该检查不同于第一阶段按字节相加的vadd补码自检。\par
\reportfigure{S2-TEST-02}{负数点积测试：低32位为−8的补码，高96位清零}{5.0cm}{八个−1与八个1点积；实际低32位为FFFFFFF8，高96位为零。}{用户保存的原始PNG；functional-0观察器日志字段摘录，保留事务身份、原始行号与SHA256；不是波形截图。}
\clearpage
\subsection{题目2.3.3：边界值点积测试}
最大正结果使用八个−128与八个−128，每项乘积16384，总和131072。functional-0的seq=5、ROB=0:43、pdest=66在448417周期接受、448419周期响应；实际完整结果为\seqsplit{0x00000000000000000000000000020000}。该例检验乘法和归约在正端点不被截断。\par
\reportfigure{S2-TEST-03}{边界值点积测试：最大正结果131072}{5.0cm}{八对−128相乘求和；实际低32位为00020000，高96位为零。}{用户保存的原始PNG；functional-0观察器日志字段摘录，保留事务身份、原始行号与SHA256；本图只证明所示最大正结果案例。}
零输入用八个0与八个127，实际完整结果为零；functional-0的seq=0、ROB=0:8、pdest=124在311084周期接受、311086周期响应。最小结果的数学预期为8×(−128)×127=−130048，低32位为\seqsplit{0xfffe0400}；本稿未单独展开该项全核输入与响应配对，因此该值只列为预期边界，不新增通过声明。\par
\noindent\textbf{证据来源与判读：}固定0f历史运行，原始记录根目录为\seqsplit{/home/riscv/riscv-lab/software-import-20261026/clean-02/fullcore-run-01/runs/}；正数在\seqsplit{functional-2/run.log}第7030—7034行，负数、最大正结果、零输入分别在\seqsplit{functional-0/run.log}第7927—7931、8400—8404、5814—5818行。按seq、ROB与pdest配对，比较完整128位响应。上述三图是日志字段摘录，不是GTKWave截图；执行响应不等同于架构提交，完整运行终态与差分结论按原归档核对。\par

\subsection{算术数据通路测试汇总}
除前述具体用例外，历史组合算术RTL测试统计为22项定向测试、10000项随机测试，合计10022项全部通过。该汇总与前述全核用例摘录属于不同验证层级，各自保留原测试范围。\par
\reportfigure{S2-TEST-SUMMARY}{VDot算术数据通路定向与随机测试结果}{4.0cm}{定向测试22项、随机测试10000项，合计10022项全部通过。本结果属于算术数据通路验证；全核协同仿真、流水线反压及取消行为由相应测试另行验证。}{用户提供的原始统计截图，完整保留三行PASS输出。本次仅补图，未重新执行测试；后续S2-09综合图中的算术统计为同一组历史结果，不重复计数。}


图\ref{fig:S2-09}把上述三类测试放在一起对照。三组输入的源高64位都填入A5哨兵，不参与运算；三个结果分别为204、$-8$和131072，写回的高96位均为零，与期望一致。底部的10022项统计与图\ref{fig:S2-TEST-SUMMARY}是同一组算术测试。


\reportfigure{S2-09}{正数、负数与边界测试的输入和结果}{12.0cm}{集中对照正数204、负数$-8$和最大正结果131072，观察低32位值及高96位零；算术统计与前述汇总为同一组数据。}{对应测试记录与结果示意，详细响应摘录见本节三张原始观察图。}

算术RTL之外，团队还用Python在主机上实现了一个独立的参考模型，用来定义指令语义并生成测试向量；它的35个单元测试也全部通过，输出见图\ref{fig:SUP-S2-MODEL}。这里按测试方法计数，与算术RTL的10022组输入分别统计。

\evidencefigure{SUP-S2-MODEL}{独立主机模型的测试输出}{11cm}{35个主机模型测试方法通过。该统计以测试方法为单位；算术RTL的10022组输入另见基础测试汇总。}

在基础数值之外，测试还检查源高半区哨兵、旧目的哨兵、寄存器别名、v0目的、非法配置和向量状态恢复。在FU和真实父级上，还向响应身份、写使能、取消状态和结果高位注入错误，确认检查器能够发现；正常运行时则核对完整响应和两槽排空。

\section{设计取舍与扩展}
\subsection{两槽结构与结果保存}\label{subsec:2-tradeoff}
采用两槽结构后，相邻请求的乘法和归约可以重叠执行，连续的独立点积因此能每拍发起一笔。将R中的结果保存宽度收窄至32位，利用了高96位固定为零这一架构性质，减少了数据寄存器成本。每笔请求的ROB、目的和控制信息则完整保留，以支持乱序执行、背压和精确取消。

运算拍数方面，团队还实现并验收了把乘法和归约合并为一拍的SP1方案。它使依赖链周期再降低约14\%，但FU面积增加约6\%，VDot从输入到寄存器的路径时序余量缩小，批量吞吐和函数调用周期基本不变。综合权衡后，最终提交仍采用两拍算术的两槽实现。两版的周期、面积和时序对比见第三阶段。

\subsection{十六项INT8点积的功能进展}\label{sec:2-dot16}
在八项实现的基础上，团队增加了独立的十六项点积候选，以补丁形式提供并附完整验证记录。它采用\texttt{MATCH=0x0e00000b}、\texttt{MASK=0xfe00707f}，在e8/m1、VL=16时读取两个128位源的全部16个字节。八项指令继续使用原编码和VL=8语义。

候选复用两组八项乘积及归约模块。P保存16个16位乘积，两组19位部分和相加得到20位精确和，再扩展为低32位结果。十六项和的范围为$[-260096,262144]$，R仍保存32位，写回高96位为零。接收八项指令时，高八个乘积写零，使后级统一完成归约。两槽握手和取消机制沿用原结构。

2026年11月6日，该候选完成了编译开关关闭和开启两种配置下的译码、FU和真实父级验证。启用十六项扩展后，使用新EMU与匹配NEMU完成六例整核差分：2例八项指令、4例十六项指令全部通过。六例覆盖源高半区贡献、有符号正负端点、选定源目的别名、完整128位写回和恢复八项指令。相应FU与父级连续请求场景还核对了2拍首响应与1拍发起间隔。

目前该扩展已完成指令和执行通路的功能验证。它每条处理16项，可以减少量化GEMV所需的点积指令条数；完整模型调用的性能对比和综合时序分析是下一步工作。本报告的周期和成本数据来自八项受测版本。

\subsection{元素位宽的参数化验证}\label{sec:2-int16}
另一组独立模块实验把输入元素扩展为INT16，计算8项点积，并使用低64位结果字段、高64位清零的输出布局。该配置的最大正和为$2^{33}$，精确累加至少需要35位。参数化乘法与归约模块完成94项定向输入和10000项随机输入的完整128位结果检查，并验证单请求包装中的反压、取消与复位行为。这组实验说明，改变元素位宽时，可以按同样的方法重新确定中间位宽、结果字段和测试向量。


% 本节标题、说明和表格放在同一页，排不下时整体移到下一页顶部。
\needspace{0.42\textheight}
\section{源码与证据入口}
表\ref{tab:s2-entries}按内容列出本章涉及的源码与证据入口。各入口下保存完整源码版本、运行输入、原始日志和失败记录，可用于复查实现细节与测量依据。
\begin{center}
\small\captionsetup{type=table}
\caption{第二阶段源码与证据入口}\label{tab:s2-entries}
\begin{tabular}{p{0.25\linewidth}p{0.64\linewidth}}
\toprule
内容 & 阅读入口\\
\midrule
指令语义与编码 & \href{https://github.com/cc77-builder/riscv/blob/a0f13456e0c3180fe9c4023f5cbe7874eb5d35a9/team/spec/vdot.md}{指令规格与定向向量}\\
算术及两槽实现 & \href{https://github.com/cc77-builder/riscv/blob/0f12df675378556670afd0dd2153f47fcb330610/src/main/scala/xiangshan/backend/fu/wrapper/VDotElastic.scala}{八项弹性FU源码}及\href{https://github.com/cc77-builder/riscv/blob/0f12df675378556670afd0dd2153f47fcb330610/src/main/scala/xiangshan/backend/fu/vector/VDotDatapath.scala}{乘积与归约源码}\\
构建与版本选择 & \href{https://github.com/cc77-builder/riscv/blob/f5f67e8640d8fda412e8a8946ffb8168f67d79fa/team/docs/development-route-current.md}{工程版本入口与可选配置}\\
周期、面积与STA & \href{https://github.com/cc77-builder/riscv/blob/a0f13456e0c3180fe9c4023f5cbe7874eb5d35a9/team/optimizations/o2/r3-div-gate/PERFORMANCE_COST_RESULT.md}{八项两槽受测版本的数据与原始报告}\\
十六项扩展 & \href{https://github.com/cc77-builder/riscv/blob/f5f67e8640d8fda412e8a8946ffb8168f67d79fa/team/optimizations/o2/dot16/rtl-integration/README.md}{候选结构、补丁与分层验证}、\href{https://github.com/cc77-builder/riscv/blob/f5f67e8640d8fda412e8a8946ffb8168f67d79fa/team/optimizations/o2/dot16/rtl-integration/fullcore/RESULT.md}{六例整核差分结果}\\
元素位宽参数化 & \href{https://github.com/cc77-builder/riscv/blob/a0f13456e0c3180fe9c4023f5cbe7874eb5d35a9/team/docs/t15-parameter-extension.md}{八项INT16独立模块及完整结果验证}\\
\bottomrule
\end{tabular}
\end{center}
